% Generated by build/make_tables.py from gs_var_average_daily.csv, gs_var_period_end.csv and var_fixed_book_daily.csv. Do not edit.

\begin{table}[tbp]
\centering
\small
\caption{Amplitude and persistence: the two disclosed measures of the diversification effect against their proxy analogues}
\label{tab:amplitude}
\begin{tabular}{@{}lrrlrr@{}}
\toprule
 & & Std. & & Change, first & Lag-1 \\
Series & Mean & dev. & Range & half to second & autocorr. \\
\midrule
Goldman, average daily & 42.03 & 2.08 & 39.0--45.5 & +0.8 & -0.47 \\
Goldman, period-end & 38.21 & 6.72 & 28.6--48.6 & +8.6 & +0.07 \\
\addlinespace
Proxy, average daily & 41.60 & 6.33 & 31.9--49.3 & -8.5 & +0.88 \\
Proxy, final day of year & 40.62 & 6.68 & 30.6--47.8 & -10.2 & +0.90 \\
\bottomrule
\end{tabular}
\notetext{All figures are the diversification effect as a percentage of summed category VaR across the ten fiscal years, so each column rests on ten observations and none is precisely estimated. \emph{Change} is the mean of 2021--2025 less the mean of 2016--2020. \emph{Lag-1 autocorr.} is the correlation between consecutive annual values: negative indicates a series that returns toward its mean, positive one that carries its level forward. Neither disclosed series shows positive persistence --- the average daily series is $-0.47$ and the period-end series is essentially zero --- while both proxy series sit near $+0.9$. That contrast, rather than any difference in spread, is what Section~\ref{sec:results} rests on. No trend test is reported: the proxy's annual values are computed from overlapping five-year windows and share about 82\% of their scenario set with their neighbours, so a regression on them would understate its own standard errors.}
\end{table}
